\section{结论}

词汇侧，家国同构确实压在极少数字上：族 234、孝 167、嗣 114、继承 82 四词合计近六百次，
而「宗族」「家法」各 1 次。通俗历史讲的是日常称谓，不是人类学概念——
把「宗族」放进这套叙事，是后人加的标签，不是书里的词。

关系侧，亲缘没有成为组织原则：亲属 7.5 对基线 8.0 持平，君臣从属 18.2 与敌对冲突 12.6
维持基线位次，未归类「其他」反而由 39.9 升到 49.1。家国同构在本书的关系结构里
主要表现为君臣框架，词汇上的「家」与关系上的「亲缘」是两套量级，不能互相折算。

四个锚点的热度与度数 $\rho=0.871$，近乎同一排布：朱元璋 41 章 / 87 度、朱棣 113 度 / 25 章。
这套语法靠的是贯穿全书的人物骨架，不是某一时期的专题。可检验的下一步是换一部年代更早、
作者不同的明代中枢叙事复算，看亲属边占比是否仍停在 8\% 上下——若仍停在 8\%，
那才是「家国同构」这个比喻在明代叙事传统里的真实份量；若抬到 20\% 以上，
那它只是这部特定通俗叙事的取景选择。

本文没有推翻什么，也没有证实什么历史新解。它做的一件小事是：把一个只能靠读感判断的母题，换成一组可以重算、可以点数、可以在另一部文本上复核的数字。这件事做成的标准只有一个——换一个人拿同一部书重跑，应当得到同样的两个小数、同一对方向。
